\documentclass[10pt,a4paper]{amsart}
\usepackage[margin=19mm]{geometry}
\usepackage{amsmath,amssymb,mathtools}
\usepackage[scr=rsfs,cal=euler]{mathalfa}
\usepackage{xfrac}
\usepackage[unicode,hidelinks,bookmarks=false]{hyperref}
\newcommand{\ie}{i.e.\ }
\newcommand{\cf}{cf.\ }
\newcommand{\resp}{respectively\ }
\newcommand{\Subsection}{\subsection}
\providecommand{\nobreakdash}{\nobreak\hbox{-}}
\setlength{\emergencystretch}{2em}
\multlinegap0pt
\usepackage{amssymb}
\usepackage{mathtools}
\usepackage{tikz}
\usepackage[shortlabels]{enumitem}
\newtheorem{lemma}{Lemma}
\newtheorem{theorem}{Theorem}
\title{A proof of Seymour's second neighborhood conjecture for oriented graphs with minimum out-degree equal to \(7\)}
\author{Arpan Sadhukhan, R. B. Sandeep, Sagnik Sen}
\date{Source: arXiv:2606.30588v1 (CC BY 4.0)}
\begin{document}
\maketitle

\noindent\emph{Source and licence.} A proof of Seymour's second neighborhood conjecture for oriented graphs with minimum out-degree equal to \(7\). Version arXiv:2606.30588v1 (CC BY 4.0), distributed by arXiv under the Creative Commons Attribution 4.0 International licence, http://creativecommons.org/licenses/by/4.0/, as recorded in the arXiv licence metadata for this exact version and shown on the abstract page https://arxiv.org/abs/2606.30588v1. Full licence notice: \texttt{licenses/arxiv-2606.30588-CC-BY-4.0.txt}.\par\smallskip

\noindent\textit{Editorial scope.} Counted unit: the article's theorem theorem (source0.tex lines 246--248). The article's complete original proof of it is delivered with it (source0.tex lines 250--326, one \texttt{proof} environment of 77 lines). No mathematics is written, repaired or re-derived: the statement and the proof are the article's own lines, in the article's own notation, and no retained line is substituted or re-flowed.  Retained uncounted as the source-local context the proof needs: lines 168--174.  Nothing was omitted from any retained span, so the package carries no omission record.  No retained line contains a citation, so the wrapper carries no bibliography and no bibliography entry.  No retained line contains a figure environment, an image-inclusion command or drawing code, so no image is delivered and the package declares no asset dependency.  Editorial formatting only, in addition: an AMS \texttt{amsart} preamble that restates the article's own declarations and macros used by the retained lines, namely the environments \texttt{lemma}, \texttt{theorem} on separate counters, together with the standard packages those lines need.  The retained environments print in this document as Lemma 1, Theorem 1 in order of appearance, whereas the article prints the same items with the numbering of its own full document.  The complete native source of the article as distributed in the arXiv e-print for this exact version is bundled unchanged as \texttt{sources/204036/source0.tex}.
\begin{lemma}~\label{lem:perfectdegree}
Suppose that \(|A|\) is odd and
\[
|A_{1}|=\frac{|A|-1}{2}.
\]
Then D[A] is a regular tournament and $|A_2|=|A_{1}|.$
\end{lemma}

\begin{theorem}
Let $D$ be a finite oriented graph with $\delta^+(D)=7$. Let $s$ be a vertex with $d^+(s)=7$, put $A=N^+(s)$ and $B=N^{++}(s)$. If $|B|=5$, then some vertex of $A$ is a Seymour vertex.
\end{theorem}

\begin{proof}
We know that every out-neighbor of a vertex of $A$ lies in $A\cup B$. Since every vertex of $D$ has outdegree at least $7$ and $|B|=5$, each vertex of $A$ has at least two out-neighbors in $A$, that is, $|N^+(a)\cap A|\ge 7-|B|=2$ for all $a\in A$.

Recall that $a_1\in A$ is the vertex of $A$ that minimizes $|N^+(a)\cap A|$ over $a\in A$, and put and $A_1=N^+(a_1)\cap A$.
Since $e(D[A])\leq \binom{7}{2}=21$, it is easy to see $|A_1|\leq 3$. Thus $|A_1|\in\{2,3\}$.

We will now handle the two cases separately.

\medskip
\noindent\emph{Case 1: $|A_1|=2$.}


Since $N^+(a_1)\subseteq A\cup B$, $|A_1|=2$, and $|B|=5$, the set $A_1\cup B$ has exactly seven vertices. Now out-degree of $a_1$ is at least seven, thus $N^+(a_1)=A_1\cup B$ or in other words $a_1$ out-dominates all vertices of $B$ and $d^+(a_1)=7$.

Recall that $A_2=(N^+(A_1)\cap A)\setminus A_1.$ By the minimality of $a_1$, each vertex of $A_1$ has at least two out-neighbors in $A$, so the two vertices of $A_1$ send at least four edges into $A$. At most one of these edges can stay inside $A_1$. Therefore at least three edges go from $A_1$ to $A\setminus (A_1\cup\{a_1\})$. A fixed vertex can receive at most two of these edges from the two vertices of $A_1$, hence $|A_2|\geq 2$.


Since $a_1$ out-dominates every vertex of $B$, every vertex of $C$ is a second out-neighbor of $a_1$, and every vertex of $A_2$ is also a second out-neighbor of $a_1$. The sets $A_2$ and $C$ are disjoint and neither is contained in $N^+(a_1)$. Thus, if $|C|\ge 5$, then,
\[
|N^{++}(a_1)|\ge |A_2|+|C|\ge 2+5=7=d^+(a_1),
\]
so $a_1$ is Seymour. Hence for the sake of contradiction assume $ |C|\le 4.$
We now count edges between the layers. Every vertex of $A$ has out-degree at least $7$ and all its out-neighbors lie in $A\cup B$, so
\[
e(A,A)+e(A,B)\ge 7\cdot 7=49.
\]
Using $e(A,A)\le 21$, we get $e(A,B)\ge 28.$ There are $7\cdot 5=35$ cross-pairs between $A$ and $B$, thus, $e(B,A)\le 35-28=7.$
The vertices of $B$ have total out-degree at least $5\cdot 7=35$, and their out-neighbors lie in $A\cup B\cup C$. Since $e(B,B)\le \binom{5}{2}=10$, We have,
$e(B,C)\ge 35-7-10=18$.

Each vertex of $C$ receives at most five edges from $B$, hence $|C|\ge \lceil 18/5\rceil=4$. Thus we may assume $C=4$.





Thus, $e(B,C)\le 5|C|=20$ and $e(B,B)\le 10$, thus by the inequality $e(B,A)+e(B,B)+e(B,C)\ge 35$, we have, $e(B,A)\ge 35-10-20=5.$
If $|A_2|\ge 3$, then clearly, $|N^{++}(a_1)|\ge |A_2|+|C|\ge 3+4=7=d^+(a_1)$,
so $a_1$ is Seymour. Thus, we may assume $|A_2|=2.$

Let $A_3=A\setminus (\{a_1\}\cup A_1\cup A_2).$ Then $|A_3|=2$. 

If some vertex of $B$ sent an edge to a vertex of $A_3$, that vertex of $A_3$ would be a second out-neighbor of $a_1$ through $B$, and then
\[
|N^{++}(a_1)|\ge |A_2|+|C|+1=2+4+1=7=d^+(a_1).
\]
Thus, if $a_1$ is not Seymour, no edge goes from $B$ to $A_3$. Also, no edge goes from $B$ to $a_1$, because $a_1$ out-dominates all vertices of $B$. Hence, every edge from $B$ to $A$ lands in $U=A_1\cup A_2.$ We also have $e(B,A)= e(B,U)\ge 5$. Now we will derive a contradiction by showing this cannot happen.


Let $m=e(A_1,A_2)$. Clearly from definition it follows that $m\geq 2$. The edges from $A_1$ to $A$ can only go inside $A_1$ or to $A_2$: they cannot go to $a_1$, and they cannot go to $A_3$ by the definition of $A_2$. Since at most one edge lies inside $A_1$, $e(A_1,A)\le m+1$.
Next count edges from $A_2$ into $A$. There are at most $2$ edges from $A_2$ to $a_1$, at most $4-m$ edges from $A_2$ to $A_1$, at most one edge inside $A_2$, and at most $|A_2||A_3|=4$ edges from $A_2$ to $A_3$. Therefore
\[
 e(A_2,A)\le 2+(4-m)+1+4=11-m.
\]
Adding together we have, $e(U,A)\le 12$.
The four vertices of $U$ have total outdegree at least $4\cdot 7=28$, and all their out-neighbors lie in $A\cup B$. Hence
\[
 e(U,B)\ge 28-e(U,A)\ge 16.
\]
Between $U$ and $B$ there are only $4\cdot 5=20$ cross-pairs. By orientation and (23),
\[
 e(B,U)\le 20-16=4,
\]
which is a contradiction. This finishes the proof for \emph{Case 1}.

\medskip
\noindent\emph{Case 2: $|A_1|=3$.}
Since $a_1$ minimizes $|N^+(a)\cap A|$ and $|A_1|=3$, by Lemma~\ref{lem:perfectdegree}, we know that $D[A]$ is a tournament and for every vertex $a\in A$, $|N^+(a) \cap A|=3$ and also $|N^{++}(a)\cap A|=3$. Now since all vertex of $A$ has the same out-degree inside $A$, by definition of $a_1$, we choose it in such a way that $|N^+(a)\cap B|$ is minimized when $a=a_1$.

Now observe that if $|N^+(a_1)\cap B|=5$ then $|N^+(a)\cap B|=5$ for all $a\in A$ which implies $e(B,A)=0$, thus $e(B,C)\geq 35-10=25$, thus $|C|\geq 5$. Also, $C\subseteq N^{++}(a_1)$ and $C\cap A=\emptyset$. Thus $|N^{++}(a_1)|=|N^{++}(a_1)\cap C| + |N^{++}(a_1)\cap A|\geq 5+3=8=|N^+(a_1)|$. Thus, $a_1$ is a Seymour vertex. Hence, we may assume $|N^{+}(a_1)\cap B|=4$ (minimum degree of $a_1$ being greater than or equal to $7$ forces the equality here).

Also, from the analysis of \emph{Case 1}, we know that $e(B,C)\geq 18$ (this was derived independent of the size of $A_1$ thus it holds in this case too). Let $E=N^{+}(a_1)\cap B$. Now, observe that, $|N^+(E)\cap C|\leq 3$, otherwise $a_1$ would be Seymour vertex. Thus, $e(E,C)\leq 12$ (as $|E|=4$). Now $e(B,C)=e(E,C)+e(B\setminus E,C)$. Thus, $e(B\setminus E,C)\geq 6$. Now, $B\setminus E$ consists of a singleton vertex, say $u$. Clearly, there exists $a^*\in A$, such that $u\in N^+(a^*)$. Thus $|N^{++}(a^*)|=|N^{++}(a^*)\cap C| + |N^{++}(a^*)\cap A|\geq 3+6=9$. Also, $|N^+(a^*)|\leq 8$. Thus, $a^*$ is a Seymour vertex. This finishes the proof of \emph{Case 2}.




\end{proof}

\end{document}
